% English translation of checked_pages/esq_119_010.tex; source: cote 119, site 10, PDF page 11, typescript page 8.
% AI-translated by GPT-6 Astra (OpenAI, ChatGPT Pro mode).
\scanpage{119/10}
\mspage{8}

from level O), one finds \uline{twelve} \uline{fundamental} ``\uline{building blocks}''
(6 of genus O, 6 of genus 1) in a ``Légo-Teichmüller game (large
box), where the marked points on the surfaces under consideration are replaced
by ``holes'' with boundary, so as to obtain surfaces with boundary,
thus building blocks that can be fitted together with gentle
friction, as in the ordinary Légo game dear to our children
(or grandchildren...). Assembling them provides a thoroughly
visual way of constructing every type of surface (it is essentially
these assemblies that will be the ``base points'' for
our famous tower), and also of visualizing the \uline{elementary} ``\uline{paths}''
through equally concrete operations, such as ``twists'', or
automorphisms of the pieces in the game, and of writing the \uline{fundamental relations}
between composite paths. Depending on the size (and the price !) of the
construction box used, one moreover \unsure{finds}\tnote{TN-010-01}{The French reading is retained as ``trcuve|trouve'' because the letterform is unresolved. Both alternatives are represented here by ``finds'', with the source doubt retained.}
many different descriptions of the Teichmüller tower by generators and
relations. The smallest box is reduced to pieces that are all
identical, of type (O,3) - these are Thurston's ``pairs of pants'',\tnote{TN-010-02}{``Pairs of pants'' translates ``pantalons'', the surface name used by the author. The letter O in the scanned type expressions is retained.}
and the Légo-Teichmüller game I am trying to describe, arising from
motivations and reflections in absolute algebraic geometry over the field $\mathbb{Q}$,
is very close to Thurston's game of ``hyperbolic geodesic surgery'',
whose existence I learned of last year from Yves Ladegaillerie.
In a microseminar with Carlos Contou-Carrère and Yves \var{Ladegraillerie}{Ladegaillerie}{slip corrected}{V-FR-010-001},
we have begun a reflection one of whose aims is to confront
the two viewpoints, which complement each other.

I add that each of the twelve building blocks in
the ``large box'' is equipped with a canonical cellular decomposition,
stable under all symmetries, whose only vertices are the
``marked points'' (or centers of the holes), and whose edges are certain
geodesic paths (for the canonical riemannian structure on the sphere
or torus under consideration) between certain pairs of vertices (namely those
that lie on the same ``real locus'', for a suitable real structure
on the complex algebraic curve under consideration). Consequently,
all the surfaces obtained in this game by assembly are equipped
with canonical cellular structures, which in their turn (cf.§3 below)

\margin{10}

